\documentclass[12pt,a4paper]{article}
\usepackage[utf8]{vietnam}
\usepackage{amsmath,amssymb}
\usepackage{tabularx}
\usepackage{tikz}
\usetikzlibrary{calc,angles,quotes,arrows.meta,patterns}
\usepackage{xcolor}
\usepackage{enumitem}
\usepackage{fancyhdr}
\usepackage[top=1.2cm, bottom=1.2cm, left=1.5cm, right=1.5cm]{geometry}
\usepackage[dethi]{ex_test}
\SAOPTN{widthSA=2.6,heightSA=0.62}

\setlength{\headheight}{15pt}
\setlength{\parskip}{0pt}

\pagestyle{fancy}
\fancyhf{}
\rhead{\textit{Môn Toán 6 -- Đề số 10}}
\lhead{\textit{Bộ đề ôn tập Giữa học kì I}}
\cfoot{\thepage}

\begin{document}

% --- KHUNG TIÊU ĐỀ ĐỀ THI ---
\noindent
\begin{tabularx}{\linewidth}{X r}
\textbf{TRƯỜNG THCS: \dotfill} & \textbf{ĐỀ KIỂM TRA ĐỊNH KỲ GIỮA HỌC KỲ I} \\
\textbf{Họ và tên: \dotfill} & \textbf{NĂM HỌC 2026 -- 2027} \\
\textbf{Lớp: \dotfill \quad SBD: \dotfill} & \textbf{Môn: TOÁN 6} \\
\hline
\end{tabularx}
\vspace{2mm}

\begin{center}
  \textbf{\large ĐỀ ÔN TẬP GIỮA HỌC KỲ I -- ĐỀ SỐ 10}\\
  \textit{Thời gian làm bài: 90 phút (Không kể thời gian phát đề)}
\end{center}
\vspace{1mm}

% --- PHẦN I ---
\noindent\textbf{\large PHẦN I. CÂU TRẮC NGHIỆM NHIỀU PHƯƠNG ÁN LỰA CHỌN (3,0 điểm)}

\textit{Thí sinh trả lời từ câu 1 đến câu 12. Mỗi câu hỏi thí sinh chỉ chọn một phương án duy nhất.}

\vspace{1mm}
\begin{ex}
Cho hai tập hợp $A = \{1; 2; 3; 4; 5\}$ và $B = \{a; b; c; d\}$. Có bao nhiêu tập hợp gồm hai phần tử thỏa mãn có một phần tử thuộc tập hợp $A$ và một phần tử thuộc tập hợp $B$?
\choice
{$16$}
{$18$}
{\True $20$}
{$9$}
\end{ex}

\begin{ex}
Tìm số tự nhiên $x$, biết $x + 20 = 42$:
\choice
{\True $22$}
{$62$}
{$2$}
{$18$}
\end{ex}

\begin{ex}
Số $308$ đọc là:
\choice
{Ba trăm tám mươi}
{Ba mươi tám}
{\True Ba trăm linh tám}
{Ba linh tám}
\end{ex}

\begin{ex}
Hợp số là:
\choice
{Số tự nhiên có nhiều hơn 2 ước}
{Số tự nhiên có 4 ước}
{\True Số tự nhiên lớn hơn 1, có nhiều hơn 2 ước}
{Số tự nhiên lớn hơn 1, chỉ có hai ước là 1 và chính nó}
\end{ex}

\begin{ex}
Trong tam giác đều, mỗi góc có số đo bằng:
\choice
{\True $60^\circ$}
{$45^\circ$}
{$90^\circ$}
{$30^\circ$}
\end{ex}

\begin{ex}
Khẳng định nào dưới đây là đúng?
\choice
{Hình bình hành có 4 góc bằng nhau}
{\True Hình lục giác đều có 3 đường chéo chính cùng đi qua một điểm}
{Tam giác đều là tam giác có 3 góc vuông}
{Hình chữ nhật có hai đường chéo bằng nhau và vuông góc với nhau}
\end{ex}

\begin{ex}
Kết quả của phép tính $7^{12} \cdot 7^2 \cdot 7^0$ dưới dạng một lũy thừa là:
\choice
{\True $7^{14}$}
{$7^{24}$}
{$7^{32}$}
{$7^7$}
\end{ex}

\begin{ex}
Cho hình bình hành $ABCD$ có $AB = 3\text{ cm}, BC = 2\text{ cm}$. Chu vi hình bình hành $ABCD$ là:
\choice
{$6\text{ cm}$}
{\True $10\text{ cm}$}
{$12\text{ cm}$}
{$5\text{ cm}$}
\end{ex}

\begin{ex}
Tập hợp nào sau đây là tập hợp các số nguyên $\mathbb{Z}$?
\choice
{$\{1; 2; 3; 4; 5\}$}
{\True $\{\dots; -3; -2; -1; 0; 1; 2; 3; \dots\}$}
{$\{0; 1; 2; 3; 4; 5\}$}
{$\{0; -1; -2; -3; \dots\}$}
\end{ex}

\begin{ex}
Số 18 được viết bằng chữ số La Mã là:
\choice
{XII}
{XV}
{\True XVIII}
{XX}
\end{ex}

\begin{ex}
Số nào sau đây là bội của $17$?
\choice
{\True $51$}
{$54$}
{$56$}
{$63$}
\end{ex}

\begin{ex}
Cho hình thoi có độ dài hai đường chéo là $2\text{ cm}$ và $4\text{ cm}$. Diện tích của hình thoi đó là:
\choice
{\True $4\text{ cm}^2$}
{$6\text{ cm}^2$}
{$8\text{ cm}^2$}
{$2\text{ cm}^2$}
\end{ex}


\vspace{1mm}
\noindent\textbf{BẢNG TRẢ LỜI TRẮC NGHIỆM PHẦN I:}
\begin{center}
\begin{tabular}{|c|c|c|c|c|c|c|c|c|c|c|c|c|}
\hline
\textbf{Câu} & \textbf{1} & \textbf{2} & \textbf{3} & \textbf{4} & \textbf{5} & \textbf{6} & \textbf{7} & \textbf{8} & \textbf{9} & \textbf{10} & \textbf{11} & \textbf{12} \\
\hline
\textbf{Chọn} & \rule{0pt}{0.65cm} & & & & & & & & & & & \\
\hline
\end{tabular}
\end{center}
\vspace{2mm}

\setcounter{ex}{0}
\noindent\textbf{\large PHẦN II. CÂU TRẮC NGHIỆM ĐÚNG SAI (2,0 điểm)}

\textit{Thí sinh trả lời câu 1 và câu 2. Trong mỗi ý a), b), c), d) ở mỗi câu, thí sinh chọn đúng hoặc sai.}

\vspace{1mm}
\begin{ex}
Cho tập hợp $B = \{x \in \mathbb{N} \mid 12 \le x < 18\}$. Xét tính đúng/sai của các khẳng định sau:
\choiceTF
{\True Tập hợp $B$ gồm các phần tử là $12; 13; 14; 15; 16; 17$}
{Số $18$ thuộc tập hợp $B$ (kí hiệu là $18 \in B$)}
{\True Trong tập hợp $B$, chỉ có duy nhất một số chia hết cho 5 là số $15$}
{\True Tổng của phần tử bé nhất và phần tử lớn nhất trong tập hợp $B$ là $29$}
\end{ex}

\begin{ex}
Cho hình thang cân $ABCD$ có độ dài đáy nhỏ $AB = 4\text{ cm}$, đáy lớn $CD = 10\text{ cm}$ và chiều cao $h = 5\text{ cm}$. Xét tính đúng/sai của các khẳng định sau:
\choiceTF
{\True Hai cạnh bên $AD$ và $BC$ có độ dài bằng nhau}
{\True Hai góc kề một đáy của hình thang cân $ABCD$ bằng nhau}
{Hai đường chéo $AC$ và $BD$ vuông góc với nhau}
{\True Diện tích của hình thang cân $ABCD$ là $35\text{ cm}^2$}
\end{ex}


\noindent\textbf{BẢNG TRẢ LỜI PHẦN II:}
\begin{center}
\begin{tabular}{|c|c|c|c|c||c|c|c|c|c|}
\hline
\textbf{Câu 1} & \textbf{a)} & \textbf{b)} & \textbf{c)} & \textbf{d)} & \textbf{Câu 2} & \textbf{a)} & \textbf{b)} & \textbf{c)} & \textbf{d)} \\
\hline
\textbf{Đ/S} & \rule{0pt}{0.6cm} & & & & \textbf{Đ/S} & & & & \\
\hline
\end{tabular}
\end{center}
\vspace{2mm}

\setcounter{ex}{0}
\noindent\textbf{\large PHẦN III. CÂU HỎI TRẢ LỜI NGẮN (2,0 điểm)}

\textit{Thí sinh trả lời câu 1 và câu 2. Điền kết quả số vào ô trống tương ứng.}

\vspace{1mm}
\begin{ex}
Tìm số tự nhiên $x$ sao cho $24 \ \vdots \ x$ và $10 \le x < 15$. \shortans{12}
\end{ex}

\begin{ex}
Một hình thang cân có độ dài hai đáy lần lượt là $6\text{ cm}$ và $10\text{ cm}$, chiều cao bằng $4\text{ cm}$. Tính diện tích của hình thang cân đó (theo đơn vị $\text{cm}^2$). \shortans{32}
\end{ex}


\setcounter{ex}{0}
\noindent\textbf{\large PHẦN IV. TỰ LUẬN (3,0 điểm)}

\textit{Thí sinh trình bày chi tiết lời giải vào giấy làm bài.}

\vspace{1mm}
\begin{ex} (1,0 điểm)
Thực hiện phép tính (Tính hợp lí nếu có thể):
\begin{enumerate}[label=\textbf{\alph*)}, leftmargin=0.8cm, itemsep=2pt]
\item (0,3 điểm) $6^2 : 4 + 2^3 \cdot 5$
\item (0,35 điểm) $47 \cdot 31 + 69 \cdot 47 - 200$
\item (0,35 điểm) $12 \cdot \{390 : [500 - (125 + 245)]\}$
\end{enumerate}
\end{ex}
\vspace{1.5mm}

\begin{ex} (1,0 điểm)
Tìm số tự nhiên $x$, biết:
\begin{enumerate}[label=\textbf{\alph*)}, leftmargin=0.8cm, itemsep=2pt]
\item (0,3 điểm) $x + 35 = 98$
\item (0,35 điểm) $3x - 45 = 60$
\item (0,35 điểm) $(6 - x)^3 + 11 = 19$
\end{enumerate}
\end{ex}
\vspace{1.5mm}

\begin{ex} (0,5 điểm)
Bác Minh muốn lát nền cho một căn phòng hình chữ nhật có chiều dài $16\text{ m}$, chiều rộng $5\text{ m}$. Loại gạch lát nền được sử dụng là gạch hình vuông có cạnh dài $40\text{ cm}$. Hỏi bác Minh phải sử dụng bao nhiêu viên gạch (coi mạch vữa không đáng kể)?
\end{ex}
\vspace{1.5mm}

\begin{ex} (0,5 điểm)
Chứng tỏ rằng biểu thức sau chia hết cho 21:\\$A = 1 + 4 + 4^2 + 4^3 + \dots + 4^{2024}$
\end{ex}
\vspace{1.5mm}


\begin{center}
\rule{4cm}{0.4pt}\\
\textit{--- HẾT ---}
\end{center}
\end{document}
